\sect{11. Equipo y objetos de la aventura}

Los dibujos muestran un ejemplo de cada familia; las variantes pueden compartir icono. El color de una poción o un anillo no basta para identificar su efecto.

Cada objeto tiene un lugar de uso: mano, cuerpo, inventario o mecanismo de la mazmorra. El nombre no basta para saber sus propiedades. Al coger un pergamino con el botón izquierdo se muestra el nombre de su conjuro sin consumirlo. Las pociones y varitas sin identificar conservan nombres genéricos: Leer magia permite identificarlas.

\textbf{Bonificaciones y maldiciones.} Una bonificación positiva puede mejorar el ataque, el daño o la armadura. Una maldición puede empeorar el equipo y bloquear su retirada. Guarda antes de probar objetos sospechosos. No hay una opción de campamento para eliminar maldiciones. Las criaturas que oxidan o destruyen equipo pueden hacerte perder objetos: conserva recambios.

\refspace
\subsection*{Armas y munición}

Daño base contra enemigos pequeños/medianos y grandes, respectivamente. Se expresa antes de fuerza, bonificaciones y defensas. \textbf{G}: guerrero, explorador y paladín; \textbf{M}: mago; \textbf{C}: clérigo; \textbf{L}: ladrón. Las combinaciones de clase pueden ampliar las opciones; el juego comprueba si el personaje puede usar el equipo.

\begin{reference}{Nombre}{Efecto, uso y límites}
\itemname{axe}{Hacha} & 1d8 / 1d10. Una mano; G, L. Incluye Drow Cleaver.\\[4pt]
\itemname{long-sword}{Espada larga} & 1d8 / 1d12. Una mano; G, L. Variantes: Slasher, Night Stalker, Severious y espada de adamantita.\\[4pt]
\itemname{short-sword}{Espada corta} & 1d6 / 1d8. Una mano; G, L. Slicer pertenece a esta familia.\\[4pt]
\itemname{dagger}{Daga} & 1d4 / 1d3. G, M, L. Se lanza al usarla desde la mano; recógela después. Backstabber, Flicka y Guinsoo son variantes.\\[4pt]
\itemname{spear}{Lanza} & 1d6 / 1d8. G, L. Se usa como arma arrojadiza desde la mano. No se repone si no llevas otra.\\[4pt]
\itemname{halberd}{Alabarda} & 1d10 / 2d6. Dos manos; G, L. La alabarda del jefe es una variante de esta familia.\\[4pt]
\itemname{mace}{Maza / cetro} & 1d6 + 1 / 1d6. Una mano; G, C, L. El Cetro de poder real funciona aquí como un arma de esta familia; no presupongas un poder especial de empuje.\\[4pt]
\itemname{flail}{Mangual} & 1d6 + 1 / 2d4. Una mano; G, C, L. Para golpear desde la primera fila.\\[4pt]
\itemname{staff}{Bastón} & 1d6 / 1d6. Dos manos; todas las clases. También lo puede empuñar un mago.\\[4pt]
\itemname{bow}{Arco} & 1d6 / 1d6. Dos manos; G, L. Necesita flechas en la otra mano y puede recargarlas del carcaj. Incluye el arco drow.\\[4pt]
\itemname{sling}{Honda} & 1d4 / 1d4. Dos manos; G, C, L. Usa piedras como munición. Comprueba ambas manos si aparece el aviso de falta de munición.\\[4pt]
\itemname{dart}{Dardo} & 1d3 / 1d2. G, M, L. Arrojadizo; puede reponerse desde el cinto. Los dardos especiales tienen modificadores propios.\\[4pt]
\itemname{arrow}{Flecha y piedra} & Munición. Las flechas se guardan en el carcaj y alimentan el arco. Las piedras alimentan la honda y también pueden lanzarse. Son objetos recuperables cuando quedan en el suelo.\\[4pt]
\end{reference}

\refspace
\subsection*{Armadura y protección}

Una CA más baja es mejor. Los valores base de la tabla no incluyen la destreza ni objetos mágicos. Dos objetos con nombres distintos pueden dejar la misma CA final.

\begin{reference}{Nombre}{Efecto, uso y límites}
\itemname{leather}{Armadura de cuero} & CA base 8; G, C, L. Una opción inicial para quien no pueda llevar armadura pesada.\\[4pt]
\itemname{scale}{Armadura de escamas} & CA base 6; G, C. La bonificación mágica y la destreza pueden mejorar el resultado.\\[4pt]
\itemname{chain}{Cota de malla} & CA base 5; G, C. Equípala en el cuerpo, no en una mano.\\[4pt]
\itemname{banded}{Armadura de bandas} & CA base 4; G, C. Compara el resultado final en la hoja del personaje.\\[4pt]
\itemname{plate}{Armadura de placas} & CA base 3; G, C. Incluye variantes mágicas o malditas, como la armadura de gran belleza; el nombre no garantiza una mejora.\\[4pt]
\itemname{robe}{Túnica} & Todas las clases. La túnica normal no mejora por sí sola la CA; la Túnica de defensa tiene una bonificación propia.\\[4pt]
\itemname{shield}{Escudo} & Reduce la CA en 1, más su posible bonificación. Debe estar en una mano que pueda usarse; no combina con el uso de un arma a dos manos.\\[4pt]
\itemname{bracers}{Brazales} & Su bonificación puede mejorar la CA al llevarlos en las muñecas. Compara el valor final al cambiar el equipo.\\[4pt]
\itemname{ring-protection}{Anillo de protección} & Llévalo en un dedo. Se utiliza el mejor de los dos anillos y no se suma a una armadura con bonificación mágica.\\[4pt]
\itemname{helmet}{Casco y botas} & Equípalos en las casillas de cabeza y pies. Los cascos y las botas normales no mejoran la CA.\\[4pt]
\end{reference}

\refspace
\subsection*{Herramientas, llaves y objetos especiales}

\begin{reference}{Nombre}{Efecto, uso y límites}
\itemname{spellbook}{Libro de conjuros} & El mago lo necesita para preparar magia. Ponlo en una mano y úsalo con el botón derecho para abrir los conjuros preparados. No se consume al lanzar magia.\\[4pt]
\itemname{holy-symbol}{Símbolo sagrado} & Herramienta del clérigo o paladín. En una mano abre sus oraciones. Llevarlo no sustituye la elección y recuperación de conjuros mediante el descanso.\\[4pt]
\itemname{lockpicks}{Ganzúas} & Coge las ganzúas con el puntero y úsalas sobre una cerradura adecuada. El ladrón puede intentar abrirla; un fallo no significa necesariamente que no sea una cerradura.\\[4pt]
\itemname{key}{Llaves} & Hay llaves de oro, plata, rubí, enjoyadas, de calavera, enanas y drow. Lleva la llave al hueco o cerradura correspondiente. Un color o forma parecidos no garantizan que encaje.\\[4pt]
\itemname{gem}{Gemas} & Conserva las gemas: pueden servir como ofrendas o piezas de un mecanismo. Su uso depende del lugar; no son una moneda que se gaste en una tienda.\\[4pt]
\itemname{stone-object}{Objetos de piedra} & Daga, símbolo sagrado, medallón, collar, orbe, anillo y cetro de piedra. Son piezas de los mecanismos de viaje. Su forma importa; guarda cada una y anota dónde la encontraste.\\[4pt]
\itemname{bones}{Huesos} & Restos humanos, enanos y medianos pueden estar relacionados con personajes de la aventura. Conserva los restos y presta atención a los diálogos; no se usan como comida.\\[4pt]
\itemname{letter}{Carta de encargo y notas} & Pon el documento en una mano y úsalo para leerlo. Las notas dan contexto o pistas. Son documentos, no un automapa que se actualice al caminar.\\[4pt]
\itemname{orb}{Orbe de poder} & Objeto especial de la aventura. No lo trates como una varita recargable ni como un arma ordinaria; conserva los objetos singulares hasta conocer su propósito.\\[4pt]
\itemname{necklace}{Collares y medallones} & Se llevan en el cuello. Algunos son adornos, como el Medallón de la suerte, que no da bonificaciones de combate.\\[4pt]
\end{reference}
